\paragraph{Adapting Vision--Language Models Under Domain Shift:}
Large pre-trained vision and vision--language models exhibit strong transfer across downstream tasks, but this transfer does not imply robustness under distribution shift. CLIP, SigLIP, ALIGN, and DINOv2 can degrade substantially on fine-grained recognition, medical imaging, satellite imagery, textures, and corruption benchmarks \citep{radford2021learning, zhai2023sigmoid, jia2021scaling, oquab2024dinov2, hendrycks2019benchmarking}, motivating a large adaptation literature. Prompt-learning methods (CoOp, CoCoOp, MaPLe, PromptSRC, TAP) learn task-specific context tokens while keeping backbone parameters fixed \citep{zhou2022learning, zhou2022conditional, khattak2023maple, khattak2023promptsrc, ding2025tap}. Parameter-efficient fine-tuning via LoRA modifies internal weights through low-rank updates \citep{hu2022lora}. Adapter-based and visual-prompt methods provide complementary mechanisms \citep{gao2024clipadapter, jia2022visual}, while distributionally robust optimization, source-free adaptation, and test-time prompt tuning target robustness directly \citep{sagawa2020distributionally, singha2023ad, shu2022tpt}.

However, this literature evaluates adaptation primarily through output behavior: accuracy, transfer, or calibration. These metrics answer whether adaptation improves performance, but not what it does to the internal feature basis. Whether adaptation reuses, repurposes, suppresses, or replaces pre-trained features remains unclear, yet this distinction is central for interpretability: two adapted models with similar accuracy may rely on very different internal mechanisms.

\vspace{-5pt}
\paragraph{From Output Explanations to Internal Representation Analysis:}
Attribution methods such as Grad-CAM, integrated gradients, and relevance propagation localize influential input regions but provide limited information about feature-level organization \citep{gradcam, IG, LRP}. This limitation is especially important under domain shift, where failures may arise from non-local cues such as texture, acquisition artifacts, or entangled feature reuse. Earlier representation-level work analyzed internal features through feature inversion, network dissection, and representation similarity \citep{mahendran2015understanding, bau2017networkdissection, durrani2020analyzing, csiszarik2021similarity}. For VLMs, Gandelsman et al. decompose CLIP-ViT representations into layer, head, and token contributions interpretable via text directions \citep{gandelsman2024interpreting}.

Nevertheless, architectural decompositions do not directly answer whether the underlying feature basis is disentangled. A head or direction may contribute to a prediction while combining multiple unrelated concepts. An adapted model may preserve high-level architectural pathways while changing the feature-level dictionary inside those pathways. Our work therefore targets feature-level interpretability rather than output or architectural attribution alone.

\vspace{-6pt}
\paragraph{Sparse Autoencoders as Feature-Level Interpretability Tools:}
Sparse autoencoders (SAEs) decompose neural activations into sparse, monosemantic features via overcomplete dictionaries trained to reconstruct intermediate activations under sparsity constraints \citep{elhage2022superposition, bricken2023monosemanticity, cunningham2024sparse}. Recent extensions to vision and vision--language encoders use SAE dictionaries as concept bottlenecks \citep{rao2024discover}, demonstrate causal interventions across recognition and segmentation \citep{karvonen2025saevision}, steer multimodal LLM outputs via CLIP vision-encoder SAE features \citep{pach2025monosemantic}, and explore sparse concept embeddings, steerability, and hierarchical SAE variants for visual representations \citep{bhalla2024splice, joseph2025steeringclipsvisiontransformer, zaremba2025msae, bussmann2025matryoshka}. Patch-level SAEs \citep{lim2024patchsae} are especially relevant as they operate over spatial tokens, enabling localized concept extraction that can reveal whether domain shifts manifest locally, globally, or both.

We position our work as a feature-level study of adaptation. Robustness benchmarks measure performance under shift; adaptation papers optimize that performance; text-grounded decompositions identify architectural contributions; prior SAE work extracts sparse concepts from fixed or lightly adapted models. What remains underexplored is whether the sparse feature dictionary learned from a base model continues to faithfully explain models adapted to substantially different domains. We address this gap by evaluating reconstruction fidelity, sparsity, monosemanticity, and causal faithfulness across model architectures, adaptation families, and domain shifts.